% XeLaTeX でコンパイルしてください
\documentclass[a4paper,11pt]{article}

% パッケージの読み込み
\usepackage{fontspec}
\setmainfont{Noto Sans CJK JP}[
  Script=CJK,
  Language=Japanese
]
\usepackage{amsmath, amssymb, amsthm}
\usepackage{tikz}
\usetikzlibrary{arrows.meta, positioning, calc, decorations.pathreplacing, shapes.geometric}
\usepackage{tcolorbox}
\tcbuselibrary{theorems, skins}
\usepackage{hyperref}
\usepackage{listings}
\usepackage{xcolor}

% Leanコードのスタイル設定
\lstdefinestyle{leanstyle}{
    language=ML,
    basicstyle=\small\ttfamily,
    keywordstyle=\color{blue}\bfseries,
    commentstyle=\color{green!60!black},
    stringstyle=\color{red},
    showstringspaces=false,
    breaklines=true,
    frame=single,
    numbers=left,
    numberstyle=\tiny\color{gray},
    backgroundcolor=\color{gray!10}
}

% 定理環境の設定
\theoremstyle{definition}
\newtheorem{definition}{定義}[section]
\newtheorem{theorem}[definition]{定理}
\newtheorem{proposition}[definition]{命題}
\newtheorem{lemma}[definition]{補題}
\newtheorem{example}[definition]{例}
\newtheorem{remark}[definition]{注意}
\newtheorem{corollary}[definition]{系}

% ボックス環境の設定
\newtcolorbox{keyidea}{
    enhanced,
    colback=blue!5!white,
    colframe=blue!75!black,
    title=重要なアイディア,
    fonttitle=\bfseries,
    attach boxed title to top left={yshift=-2mm, xshift=5mm},
    boxed title style={colback=blue!75!black}
}

\newtcolorbox{proofstrategy}{
    enhanced,
    colback=green!5!white,
    colframe=green!60!black,
    title=証明の戦略,
    fonttitle=\bfseries,
    attach boxed title to top left={yshift=-2mm, xshift=5mm},
    boxed title style={colback=green!60!black}
}

\newtcolorbox{geometric}{
    enhanced,
    colback=orange!5!white,
    colframe=orange!80!black,
    title=幾何学的直感,
    fonttitle=\bfseries,
    attach boxed title to top left={yshift=-2mm, xshift=5mm},
    boxed title style={colback=orange!80!black}
}

% タイトル情報
\title{構造塔の切断操作と包含射\\{\large 完全版：基礎から定理まで}}
\author{%
    Lean 形式化: CODEX (OpenAI)\\
    \LaTeX 解説: Claude (Anthropic)
}
\date{\today}

\begin{document}

\maketitle

\begin{abstract}
本稿では、Bourbaki流の構造主義的数学における「構造塔」(Structure Tower)の切断操作について、包含射の理論と補題の階層構造を含めて包括的に解説する。特に、\texttt{sliceAbove}操作の定義から始まり、包含射の性質、基礎補題、そして完全性定理に至るまでの「証明の敷き詰めパターン」を、Lean 4による形式化を通じて詳しく説明する。本解説は学部2年生以上を対象とし、直感的理解と形式的厳密さの両立を目指した。
\end{abstract}

\tableofcontents
\newpage

\section{導入と動機}

\subsection{なぜ構造塔を研究するのか}

数学では、より単純な構造から複雑な構造へと段階的に理論を構築していく。例えば：
\begin{itemize}
    \item 集合 $\to$ 群 $\to$ 環 $\to$ 体
    \item 位相空間 $\to$ 距離空間 $\to$ ノルム空間 $\to$ 内積空間
\end{itemize}

構造塔は、このような階層的な数学的構造を統一的に扱うための枠組みである。

\subsection{本稿の目標}

この文書を読み終えると、以下のことが理解できる：
\begin{enumerate}
    \item 構造塔の形式的定義
    \item 切断操作\texttt{sliceAbove}の幾何学的・代数的意味
    \item 包含射の構成と性質
    \item 証明を階層的に構築する「敷き詰めパターン」
    \item Lean 4による形式証明の書き方
\end{enumerate}

\section{構造塔の基礎理論}

\subsection{構造塔の直感的イメージ}

\begin{center}
\begin{tikzpicture}[scale=1.2]
    % 塔の描画
    \foreach \i/\name/\color in {0/集合/blue!20, 1/群/green!20, 2/環/orange!20, 3/体/red!20} {
        \fill[\color] (-2,\i*0.9) rectangle (2,\i*0.9+0.8);
        \draw[thick] (-2,\i*0.9) rectangle (2,\i*0.9+0.8);
        \node at (0,\i*0.9+0.4) {\name};
        \node[left] at (-2.3,\i*0.9+0.4) {$i_{\i}$};
    }
    
    % 要素の表示
    \foreach \i in {0,1,2,3} {
        \foreach \j in {0,1,2} {
            \fill[black] (2.5+\j*0.5,\i*0.9+0.4) circle (0.05);
        }
    }
    
    % 矢印と説明
    \draw[-Stealth, ultra thick, blue] (4,3) -- (4,1);
    \node[right, align=left, blue] at (4.2,2) {より豊かな\\構造へ};
    
    % 要素の動き
    \draw[-Stealth, thick, red] (2.5,0.4) to[out=45, in=-45] (2.5,1.3);
    \draw[-Stealth, thick, red] (3,0.4) to[out=45, in=-45] (3,2.2);
    \node[right, red, font=\small] at (3.7,1.5) {単調性};
\end{tikzpicture}
\end{center}

\subsection{形式的定義}

\begin{definition}[構造塔 with Minimum Layer]
構造塔 $T$ は以下のデータと公理からなる：

\textbf{データ：}
\begin{enumerate}
    \item \textbf{台集合} $T.\text{carrier}$ : すべての要素を含む型
    \item \textbf{添字集合} $T.\text{Index}$ : 前順序が入った型
    \item \textbf{層関数} $T.\text{layer} : T.\text{Index} \to \text{Set}(T.\text{carrier})$
    \item \textbf{最小層関数} $T.\text{minLayer} : T.\text{carrier} \to T.\text{Index}$
\end{enumerate}

\textbf{公理：}
\begin{enumerate}
    \item \textbf{被覆性}: $\forall x \in T.\text{carrier}, \exists i \in T.\text{Index}, x \in T.\text{layer}(i)$
    \item \textbf{単調性}: $i \leq j \Rightarrow T.\text{layer}(i) \subseteq T.\text{layer}(j)$
    \item \textbf{最小層の帰属}: $\forall x, x \in T.\text{layer}(T.\text{minLayer}(x))$
    \item \textbf{最小性}: $\forall x, i, (x \in T.\text{layer}(i)) \Rightarrow T.\text{minLayer}(x) \leq i$
\end{enumerate}
\end{definition}

\begin{keyidea}
構造塔の核心は\textbf{単調性}にある。下の層に属する要素は必ず上の層にも属する。これは「一度満たした性質は失われない」という数学的直感を形式化したものである。
\end{keyidea}

\section{切断操作: sliceAbove}

\subsection{幾何学的動機}

構造塔全体ではなく、「ある添字$i$より上の部分だけ」を取り出したい場合がある。

\begin{center}
\begin{tikzpicture}[scale=1.2]
    % 元の塔
    \begin{scope}[xshift=-4.5cm]
        \node[above] at (0,5.5) {\textbf{元の構造塔 $T$}};
        \foreach \i/\name in {0/層0, 1/層1, 2/層2, 3/層3, 4/層4} {
            \fill[blue!20] (-1.5,\i*0.9) rectangle (1.5,\i*0.9+0.7);
            \draw[thick] (-1.5,\i*0.9) rectangle (1.5,\i*0.9+0.7);
            \node at (0,\i*0.9+0.35) {\small \name};
        }
        
        % 要素の配置
        \foreach \i in {0,...,4} {
            \node[circle, fill=black, inner sep=1.5pt] at (-0.8,\i*0.9+0.35) {};
            \node[circle, fill=black, inner sep=1.5pt] at (0,\i*0.9+0.35) {};
            \node[circle, fill=black, inner sep=1.5pt] at (0.8,\i*0.9+0.35) {};
        }
        
        % 切断線
        \draw[red, ultra thick, dashed] (-2,2*0.9-0.1) -- (2,2*0.9-0.1);
        \node[red, left] at (-2,2*0.9-0.1) {切断 $i$};
    \end{scope}
    
    % 矢印
    \draw[-Stealth, ultra thick, blue] (-1.5,2.5) -- (1.5,2.5);
    \node[above, blue] at (0,2.7) {\texttt{sliceAbove}$(T, i)$};
    
    % 切断後の塔
    \begin{scope}[xshift=4.5cm]
        \node[above] at (0,5.5) {\textbf{切断後 $T'$}};
        \foreach \i/\name in {2/層2, 3/層3, 4/層4} {
            \fill[green!20] (-1.5,\i*0.9) rectangle (1.5,\i*0.9+0.7);
            \draw[thick] (-1.5,\i*0.9) rectangle (1.5,\i*0.9+0.7);
            \node at (0,\i*0.9+0.35) {\small \name};
        }
        
        % 残った要素
        \foreach \i in {2,3,4} {
            \node[circle, fill=green!70!black, inner sep=1.5pt] at (-0.8,\i*0.9+0.35) {};
            \node[circle, fill=green!70!black, inner sep=1.5pt] at (0,\i*0.9+0.35) {};
            \node[circle, fill=green!70!black, inner sep=1.5pt] at (0.8,\i*0.9+0.35) {};
        }
        
        % 削除された部分
        \foreach \i in {0,1} {
            \draw[dashed, gray] (-1.5,\i*0.9) rectangle (1.5,\i*0.9+0.7);
            \node[gray, font=\tiny] at (0,\i*0.9+0.35) {削除};
        }
    \end{scope}
\end{tikzpicture}
\end{center}

\subsection{形式的定義}

\begin{definition}[\texttt{sliceAbove}の定義]
構造塔$T$と添字$i \in T.\text{Index}$に対して、$\texttt{sliceAbove}(T, i)$は以下のように定義される新しい構造塔$T'$：

\begin{align*}
T'.\text{carrier} &:= \{x \in T.\text{carrier} \mid i \leq T.\text{minLayer}(x)\} \\
T'.\text{Index} &:= \{j \in T.\text{Index} \mid i \leq j\} \\
T'.\text{layer}(j) &:= \{x \in T'.\text{carrier} \mid x \in T.\text{layer}(j)\} \\
T'.\text{minLayer}(x) &:= T.\text{minLayer}(x)
\end{align*}

ここで、$T'.\text{carrier}$と$T'.\text{Index}$は\textbf{部分型}(subtype)として定義される。
\end{definition}

\begin{remark}[Leanにおける部分型]
Lean 4では：
\begin{itemize}
    \item $\{x : A \mid P(x)\}$ は集合（Set A型）
    \item $\{x : A \mathbin{//} P(x)\}$ は部分型（Type）
    \item 構造体のフィールドには\textbf{Type}が必要
\end{itemize}
したがって、\texttt{carrier}と\texttt{Index}には部分型を使用する。
\end{remark}

\subsection{切断の性質の証明}

\texttt{sliceAbove}が再び構造塔の公理を満たすことを証明する必要がある。

\begin{proofstrategy}
各公理を順に確認：
\begin{enumerate}
    \item \textbf{被覆性}: 元の塔の被覆性と最小層の性質から導く
    \item \textbf{単調性}: 元の塔の単調性を継承
    \item \textbf{最小層の帰属・最小性}: 元の塔の性質から直接従う
\end{enumerate}
\end{proofstrategy}

\begin{proof}[被覆性の証明スケッチ]
$x \in T'.\text{carrier}$とする。すると$i \leq T.\text{minLayer}(x)$。

元の塔$T$の被覆性により、ある$k \in T.\text{Index}$が存在して$x \in T.\text{layer}(k)$。

最小層の最小性により$T.\text{minLayer}(x) \leq k$。

推移性より$i \leq k$なので、$k \in T'.\text{Index}$。

したがって$x \in T'.\text{layer}(k)$が成り立つ。
\end{proof}

\section{包含射: sliceAbove\_inclusion}

\subsection{包含射の定義}

切断された塔$T'$から元の塔$T$への自然な「包含」写像が存在する。

\begin{definition}[包含射]
構造塔の準同型$f : T_1 \to T_2$は以下のデータからなる：
\begin{enumerate}
    \item \textbf{台写像} $f.\text{map} : T_1.\text{carrier} \to T_2.\text{carrier}$
    \item \textbf{添字写像} $f.\text{indexMap} : T_1.\text{Index} \to T_2.\text{Index}$
\end{enumerate}
で、以下を満たす：
\begin{enumerate}
    \item \textbf{添字単調性}: $i \leq j \Rightarrow f.\text{indexMap}(i) \leq f.\text{indexMap}(j)$
    \item \textbf{層保存性}: $x \in T_1.\text{layer}(i) \Rightarrow f.\text{map}(x) \in T_2.\text{layer}(f.\text{indexMap}(i))$
    \item \textbf{最小層保存}: $f.\text{indexMap}(T_1.\text{minLayer}(x)) = T_2.\text{minLayer}(f.\text{map}(x))$
\end{enumerate}
\end{definition}

\begin{proposition}[\texttt{sliceAbove\_inclusion}]
$T' = \texttt{sliceAbove}(T, i)$に対して、包含射
\[
\iota : T' \to T
\]
が次のように定義される：
\begin{align*}
\iota.\text{map} &:= \text{Subtype.val} : \{x \mid i \leq T.\text{minLayer}(x)\} \to T.\text{carrier} \\
\iota.\text{indexMap} &:= \text{Subtype.val} : \{j \mid i \leq j\} \to T.\text{Index}
\end{align*}
これは構造塔の準同型である。
\end{proposition}

\begin{center}
\begin{tikzpicture}[scale=1.3]
    % T' (左側)
    \begin{scope}[xshift=-3cm]
        \node[above, font=\bfseries] at (0,3.5) {$T' = \texttt{sliceAbove}(T, i)$};
        \draw[thick, fill=green!10] (-1.5,0) rectangle (1.5,3);
        \node at (0,2.5) {層4};
        \node at (0,1.5) {層3};
        \node at (0,0.5) {層2};
        
        % 要素
        \node[circle, fill=green!70!black, inner sep=2pt, label=right:$x'$] (x1) at (0.3,1) {};
        \node[circle, fill=green!70!black, inner sep=2pt, label=right:$y'$] (y1) at (-0.5,2) {};
    \end{scope}
    
    % 矢印（包含射）
    \draw[-Stealth, ultra thick, blue] (-1,1.5) -- (1,1.5) 
        node[midway, above, font=\small] {$\iota = \texttt{sliceAbove\_inclusion}$};
    
    % T (右側)
    \begin{scope}[xshift=3cm]
        \node[above, font=\bfseries] at (0,4.5) {$T$};
        \draw[thick, fill=blue!10] (-1.5,-1) rectangle (1.5,4);
        \node at (0,3.5) {層4};
        \node at (0,2.5) {層3};
        \node at (0,1.5) {層2};
        \node[gray] at (0,0.5) {層1};
        \node[gray] at (0,-0.5) {層0};
        
        % 要素
        \node[circle, fill=blue!70!black, inner sep=2pt, label=right:$x$] (x2) at (0.3,1) {};
        \node[circle, fill=blue!70!black, inner sep=2pt, label=right:$y$] (y2) at (-0.5,2) {};
    \end{scope}
    
    % 対応の矢印
    \draw[->, thick, dashed, red] (x1) to[out=0, in=180] (x2);
    \draw[->, thick, dashed, red] (y1) to[out=0, in=180] (y2);
\end{tikzpicture}
\end{center}

\subsection{包含射の性質}

\begin{theorem}[包含射は単射]
包含射$\iota = \texttt{sliceAbove\_inclusion}(T, i)$について、
\[
\iota.\text{map} : T'.\text{carrier} \to T.\text{carrier}
\]
は単射である。
\end{theorem}

\begin{proof}
$\iota.\text{map} = \text{Subtype.val}$であり、\texttt{Subtype.val}は常に単射である。
\end{proof}

\begin{keyidea}
包含射の単射性は、切断操作が「要素を捨てる」ことはあっても「要素を潰す」ことはないことを意味する。
\end{keyidea}

\section{証明の敷き詰めパターン}

\subsection{階層的証明構造}

大きな定理を証明するには、小さな補題を積み重ねる「敷き詰め」戦略が有効である。

\begin{center}
\begin{tikzpicture}[scale=1.2]
    % 屋根（定理）
    \node[draw, fill=red!20, minimum width=8cm, minimum height=1cm, 
          font=\bfseries] (theorem) at (0,4) 
          {定理: \texttt{sliceAbove\_complete}};
    
    % 第2層（中間補題）
    \node[draw, fill=orange!20, minimum width=7cm, minimum height=0.8cm] (layer2) at (0,2.8) 
          {補題3: \texttt{inclusion\_layer\_comm}};
    
    % 第1層（基礎補題）
    \node[draw, fill=yellow!30, minimum width=3cm, minimum height=0.7cm] (layer1a) at (-2,1.5) 
          {補題1: \texttt{mem\_iff}};
    \node[draw, fill=yellow!30, minimum width=3cm, minimum height=0.7cm] (layer1b) at (2,1.5) 
          {補題2: \texttt{layer\_subset}};
    
    % 土台（定義）
    \node[draw, fill=green!20, minimum width=8cm, minimum height=1cm,
          font=\bfseries] (base) at (0,0) 
          {土台: \texttt{sliceAbove}定義 + 包含射};
    
    % 矢印
    \draw[-Stealth, thick] (base.north) -- (layer1a.south);
    \draw[-Stealth, thick] (base.north) -- (layer1b.south);
    \draw[-Stealth, thick] (layer1a.north) -- (layer2.south west);
    \draw[-Stealth, thick] (layer1b.north) -- (layer2.south east);
    \draw[-Stealth, thick] (layer2.north) -- (theorem.south);
    
    % ラベル
    \node[left, font=\small] at (-4.5,0) {基礎定義};
    \node[left, font=\small] at (-4.5,1.5) {基礎層};
    \node[left, font=\small] at (-4.5,2.8) {中間層};
    \node[left, font=\small] at (-4.5,4) {屋根};
\end{tikzpicture}
\end{center}

\subsection{基礎層の補題}

\begin{lemma}[補題1: \texttt{sliceAbove\_mem\_iff}]
$T' = \texttt{sliceAbove}(T, i)$とし、$x \in T.\text{carrier}$とする。このとき：
\[
i \leq T.\text{minLayer}(x) \iff \exists h : i \leq T.\text{minLayer}(x), \langle x, h \rangle \in T'.\text{layer}\langle T.\text{minLayer}(x), h \rangle
\]
\end{lemma}

\begin{geometric}
この補題は「$x$が切断後の塔に含まれる $\iff$ $x$の最小層が切断位置より上」を述べている。
\end{geometric}

\begin{lemma}[補題2: \texttt{sliceAbove\_layer\_subset}]
$i \leq j$ならば、
\[
\{x \in T'.\text{carrier} \mid x \in T.\text{layer}(j)\} \subseteq T'.\text{layer}\langle j, h_{ij} \rangle
\]
ただし$h_{ij} : i \leq j$。
\end{lemma}

\subsection{中間層の補題}

\begin{lemma}[補題3: \texttt{sliceAbove\_inclusion\_layer\_comm}]
包含射$\iota = \texttt{sliceAbove\_inclusion}(T, i)$と$j \in T'.\text{Index}$について：
\[
\forall x \in T'.\text{carrier}, \quad x \in T'.\text{layer}(j) \Rightarrow \iota(x) \in T.\text{layer}(j)
\]
\end{lemma}

この補題は包含射が層構造と可換であることを示す：

\begin{center}
\begin{tikzpicture}[scale=1.5]
    \node (TL) at (0,2) {$x \in T'.\text{layer}(j)$};
    \node (TR) at (4,2) {$\iota(x) \in T.\text{layer}(j)$};
    \node (BL) at (0,0) {$x \in T'.\text{carrier}$};
    \node (BR) at (4,0) {$\iota(x) \in T.\text{carrier}$};
    
    \draw[-Stealth, thick] (BL) -- (TL) node[midway, left] {層$j$に属する};
    \draw[-Stealth, thick] (BR) -- (TR) node[midway, right] {層$j$に属する};
    \draw[-Stealth, thick] (BL) -- (BR) node[midway, below] {$\iota$};
    \draw[-Stealth, thick] (TL) -- (TR) node[midway, above] {$\iota$};
    
    \node at (2,-0.8) {可換図式};
\end{tikzpicture}
\end{center}

\subsection{屋根：完全性定理}

\begin{theorem}[\texttt{sliceAbove\_complete}]
$T' = \texttt{sliceAbove}(T, i)$とする。任意の$x \in T.\text{carrier}$で$i \leq T.\text{minLayer}(x)$を満たすものについて、
\[
\exists y \in T'.\text{carrier}, \quad y = x \quad (\text{via } \iota)
\]
が成り立つ。
\end{theorem}

\begin{proof}
$y := \langle x, h \rangle \in T'.\text{carrier}$とおく（ただし$h : i \leq T.\text{minLayer}(x)$）。

すると$\iota(y) = \text{Subtype.val}(\langle x, h \rangle) = x$。
\end{proof}

\begin{keyidea}
完全性定理は「元の塔で切断位置より上にいた要素はすべて切断後の塔に残る」という直感を形式化したものである。
\end{keyidea}

\section{具体例で理解する}

\subsection{代数系の階層での切断}

\begin{example}[環以上の構造を取り出す]
次のような代数系の構造塔を考える：
\begin{align*}
i_0 &: \text{集合} \\
i_1 &: \text{マグマ}(二項演算) \\
i_2 &: \text{半群}(結合律) \\
i_3 &: \text{モノイド}(単位元) \\
i_4 &: \text{群}(逆元) \\
i_5 &: \text{可換群}(交換律) \\
i_6 &: \text{環}(加法群+分配的乗法) \\
i_7 &: \text{可換環} \\
i_8 &: \text{整域} \\
i_9 &: \text{体}
\end{align*}

$\texttt{sliceAbove}(T, i_6)$を適用すると：
\begin{itemize}
    \item 残る構造：環、可換環、整域、体
    \item 消える構造：集合、マグマ、半群、モノイド、群、可換群
    \item 包含射：残った構造を元の塔に埋め込む
\end{itemize}
\end{example}

\begin{center}
\begin{tikzpicture}[scale=0.85]
    % 元の塔
    \begin{scope}[xshift=-5.5cm]
        \node[above, font=\small\bfseries] at (0,7) {元の塔 $T$};
        \foreach \i/\name/\col in {
            0/集合/blue!10,
            1/マグマ/blue!15,
            2/半群/blue!20,
            3/モノイド/blue!25,
            4/群/blue!30,
            5/可換群/blue!35,
            6/環/green!30,
            7/可換環/green!40,
            8/整域/green!50,
            9/体/green!60
        } {
            \fill[\col] (-1.3,\i*0.6) rectangle (1.3,\i*0.6+0.5);
            \draw[thick] (-1.3,\i*0.6) rectangle (1.3,\i*0.6+0.5);
            \node[font=\tiny] at (0,\i*0.6+0.25) {\name};
        }
        \draw[red, ultra thick, dashed] (-1.6,6*0.6-0.05) -- (1.6,6*0.6-0.05);
        \node[red, right, font=\tiny] at (1.6,6*0.6-0.05) {切断 $i_6$};
    \end{scope}
    
    % 矢印
    \draw[-Stealth, ultra thick, blue] (-3.5,3) -- (-1.5,3);
    \node[above, font=\small] at (-2.5,3.2) {\texttt{sliceAbove}};
    
    % 切断後の塔
    \begin{scope}[xshift=1cm]
        \node[above, font=\small\bfseries] at (0,7) {切断後 $T'$};
        \foreach \i/\name/\col in {
            6/環/green!30,
            7/可換環/green!40,
            8/整域/green!50,
            9/体/green!60
        } {
            \fill[\col] (-1.3,\i*0.6) rectangle (1.3,\i*0.6+0.5);
            \draw[thick] (-1.3,\i*0.6) rectangle (1.3,\i*0.6+0.5);
            \node[font=\tiny] at (0,\i*0.6+0.25) {\name};
        }
        % 削除された部分
        \foreach \i in {0,...,5} {
            \draw[dashed, gray, opacity=0.3] (-1.3,\i*0.6) rectangle (1.3,\i*0.6+0.5);
        }
    \end{scope}
    
    % 包含射の矢印
    \draw[-Stealth, thick, red, dashed] (2,4.5) to[out=0, in=180] (4,4.5);
    \node[right, font=\small, red] at (4,4.5) {包含射 $\iota$};
    
    % 元の塔（右側、包含先）
    \begin{scope}[xshift=6.5cm]
        \node[above, font=\small\bfseries] at (0,7) {元の塔 $T$（包含先）};
        \foreach \i/\name/\col in {
            0/集合/blue!10,
            1/マグマ/blue!15,
            2/半群/blue!20,
            3/モノイド/blue!25,
            4/群/blue!30,
            5/可換群/blue!35,
            6/環/green!30,
            7/可換環/green!40,
            8/整域/green!50,
            9/体/green!60
        } {
            \fill[\col] (-1.3,\i*0.6) rectangle (1.3,\i*0.6+0.5);
            \draw[thick] (-1.3,\i*0.6) rectangle (1.3,\i*0.6+0.5);
            \node[font=\tiny] at (0,\i*0.6+0.25) {\name};
        }
    \end{scope}
\end{tikzpicture}
\end{center}

\subsection{完全性定理の意味}

具体的に、$\mathbb{Z}$(整数環)を考える：
\begin{itemize}
    \item $\mathbb{Z}$の最小層は$i_6$(環)
    \item $i_6 \leq i_6$なので、$\mathbb{Z} \in T'.\text{carrier}$
    \item 完全性定理により、$\mathbb{Z}$は切断後の塔に含まれる
    \item 包含射を通じて、元の塔の$\mathbb{Z}$と同一視される
\end{itemize}

\section{Lean 4 による形式化の詳細}

\subsection{型理論的側面}

Leanでは部分型を明示的に扱う：

\begin{verbatim}
carrier := {x : T.carrier // i ≤ T.minLayer x}
\end{verbatim}

ここで\texttt{//}は部分型構成子で、型\texttt{T.carrier}に条件\texttt{i ≤ T.minLayer x}を満たす要素の型を作る。

部分型の要素は\texttt{⟨value, proof⟩}のペアである：
\begin{itemize}
    \item \texttt{value}: 実際の値
    \item \texttt{proof}: 条件を満たす証明
\end{itemize}

\subsection{証明の戦術}

主要な証明戦術の使用例：

\begin{enumerate}
    \item \textbf{intro}: 全称量化子や含意を導入
    \begin{verbatim}
    intro x    -- ∀ x を x に具体化
    \end{verbatim}
    
    \item \textbf{obtain}: 存在量化子を分解
    \begin{verbatim}
    obtain ⟨k, hk⟩ := T.covering x.val
    -- ∃ k, ... を k と hk に分解
    \end{verbatim}
    
    \item \textbf{calc}: 等式・不等式の連鎖
    \begin{verbatim}
    calc
      i ≤ T.minLayer x.val := x.property
      _ ≤ k := T.minLayer_minimal x.val k hk
    \end{verbatim}
    
    \item \textbf{exact}: 証明項を直接与える
    \begin{verbatim}
    exact hk   -- hk がまさに求める証明
    \end{verbatim}
    
    \item \textbf{use}: 存在量化子の witness を与える
    \begin{verbatim}
    use ⟨k, hik⟩   -- ∃ ... を示すために ⟨k, hik⟩ を提示
    \end{verbatim}
\end{enumerate}

\subsection{包含射の定義（Lean コード）}

\begin{verbatim}
def sliceAbove_inclusion (T : StructureTowerWithMin) 
    (i : T.Index) : sliceAbove T i ⟶ T where
  map := Subtype.val
  indexMap := Subtype.val
  indexMap_mono := fun hjj => hjj
  layer_preserving := fun _ _ hx => hx
  minLayer_preserving := fun _ => rfl
\end{verbatim}

ここで：
\begin{itemize}
    \item \texttt{Subtype.val}: 部分型から元の型への射影
    \item \texttt{indexMap\_mono}: 添字写像の単調性（恒等的に成立）
    \item \texttt{layer\_preserving}: 層の保存（定義から直接）
    \item \texttt{minLayer\_preserving}: 最小層の保存（反射律）
\end{itemize}

\section{まとめと今後の展望}

\subsection{本稿で学んだこと}

\begin{enumerate}
    \item \textbf{構造塔の理論}
    \begin{itemize}
        \item 階層的数学構造の統一的枠組み
        \item 単調性と最小層の重要性
    \end{itemize}
    
    \item \textbf{切断操作}
    \begin{itemize}
        \item 幾何学的直感：「塔を水平に切る」
        \item 形式的定義：部分型による制限
    \end{itemize}
    
    \item \textbf{包含射}
    \begin{itemize}
        \item 切断後の塔から元の塔への自然な埋め込み
        \item 単射性と構造保存
    \end{itemize}
    
    \item \textbf{証明の敷き詰めパターン}
    \begin{itemize}
        \item 基礎補題 $\to$ 中間補題 $\to$ 定理
        \item 階層的証明構築の重要性
    \end{itemize}
    
    \item \textbf{形式証明の技法}
    \begin{itemize}
        \item Lean 4での部分型の扱い
        \item 証明戦術の効果的な使用
    \end{itemize}
\end{enumerate}

\subsection{さらなる発展}

\begin{enumerate}
    \item \textbf{双対操作}: \texttt{sliceBelow}（添字$i$以下の部分を取る）
    
    \item \textbf{圏論的視点}: 構造塔の圏、関手としての切断操作
    
    \item \textbf{極限と余極限}: 構造塔の族からの極限構成
    
    \item \textbf{具体的応用}: 
    \begin{itemize}
        \item 代数系の理論
        \item 位相空間の階層
        \item 論理体系の層構造
    \end{itemize}
\end{enumerate}

\subsection{最後に}

構造塔の理論は、Bourbakiが提唱した「構造としての数学」の現代的形式化である。本稿で見たように、Lean 4のような証明支援系を使うことで、直感的なアイディアを厳密に形式化し、機械的に検証することができる。

数学の「建築」において、証明を敷き詰めのように積み上げる手法は、大きな定理を達成するための実践的な戦略である。基礎から屋根まで、一つ一つの補題がしっかりと支え合う構造を作ることが、美しく堅固な数学の鍵となる。

\begin{thebibliography}{9}
\bibitem{bourbaki}
N. Bourbaki, \textit{Theory of Sets}, Springer, 2004.

\bibitem{lean}
The Lean Community, \textit{Theorem Proving in Lean 4}, 
\url{https://leanprover.github.io/theorem_proving_in_lean4/}

\bibitem{mathlib}
The mathlib Community, \textit{The Lean Mathematical Library},
\url{https://github.com/leanprover-community/mathlib4}

\bibitem{typetheory}
The Univalent Foundations Program, \textit{Homotopy Type Theory: 
Univalent Foundations of Mathematics}, 
\url{https://homotopytypetheory.org/book/}
\end{thebibliography}

\appendix
\section{完全なLeanコード}

\subsection{sliceAbove の定義}

\begin{lstlisting}[style=leanstyle, caption=sliceAbove の定義]
def sliceAbove (T : StructureTowerWithMin.{u, v}) 
    (i : T.Index) : StructureTowerWithMin.{u, v} where
  carrier := {x : T.carrier // i <= T.minLayer x}
  Index := {j : T.Index // i <= j}
  indexPreorder := inferInstance
  layer j := {x | x.val in T.layer j.val}
  covering := by
    intro x
    obtain (k, hk) := T.covering x.val
    have hik : i <= k := by calc ...
    use (k, hik)
    exact hk
  monotone := by
    intro j1 j2 hjj
    intro x hx
    exact T.monotone hjj hx
  minLayer := fun x => (T.minLayer x.val, x.property)
  minLayer_mem := by intro x; exact T.minLayer_mem x.val
  minLayer_minimal := by
    intro x j hx
    exact T.minLayer_minimal x.val j.val hx
\end{lstlisting}

\subsection{包含射と完全性定理}

\begin{lstlisting}[style=leanstyle, caption=包含射と完全性]
def sliceAbove_inclusion (T : StructureTowerWithMin) 
    (i : T.Index) : sliceAbove T i -> T where
  map := Subtype.val
  indexMap := Subtype.val
  indexMap_mono := fun hjj => hjj
  layer_preserving := fun _ _ hx => hx
  minLayer_preserving := fun _ => rfl

theorem sliceAbove_complete (T : StructureTowerWithMin) 
    (i : T.Index) (x : T.carrier) 
    (hx : i <= T.minLayer x) :
    exists y : (sliceAbove T i).carrier, y.val = x :=
  ((x, hx), rfl)
\end{lstlisting}

\end{document}
